\uuid{SFsi}
\exo7id{788}
\titre{exo7 788}
\auteur{gourio}
\organisation{exo7}
\datecreate{2001-09-01}
\isIndication{false}
\isCorrection{false}
\chapitre{Calcul d'intégrales}
\sousChapitre{Théorie}
\cpge{mpsi-mp2i.integration.c, pcsi.integration.b, ptsi.integration.b, mpsi-mp2i.fonctions-valeurs-reelles-complexes.b, pcsi.techniques-calcul-fonctions.b, ptsi.techniques-calcul-fonctions.b, bcpst1.integration-segment.a, bcpst1.integration-segment.b}
\module{Analyse}
\niveau{L1}
\difficulte{}

\contenu{
\question{
Soit $f\in C([0,1],{\Rr})$ telle que :
$$\forall g\in E\left( [0,1],{\Rr}\right) ,\int_{0}^{1}fg=0. $$
Montrer que $f=0.$
}
}
